% =====================================================================
% icu_sedation_paper.tex
% =====================================================================
%
% Manuscript: A Hybrid Mechanistic-Machine Learning Framework for ICU
%             Sedation Prediction Through Patient-Specific Digital Twins
%
% Author: Christopher Morris
% Affiliation: University of Florida, Department of Computer and
%              Information Science and Engineering
%
% Target journals (in order of preference):
%   1. NPJ Digital Medicine
%   2. British Journal of Anaesthesia
%   3. Critical Care Medicine
%   4. Computers in Biology and Medicine
%
% This is the working draft of the paper manuscript. Each section
% contains placeholders in italics indicating where final content
% needs to be filled in once experiments are complete. Comments
% throughout the file explain what each section should accomplish
% and provide guidance on writing decisions.
%
% =====================================================================

% Document class declaration. We use the article class for now since
% it is the most flexible. When we choose a specific target journal
% we will replace this with the journal's template.
\documentclass[11pt,a4paper]{article}

% =====================================================================
% PREAMBLE: Package loading and document configuration
% =====================================================================

% Standard packages for English language documents
\usepackage[utf8]{inputenc}
\usepackage[T1]{fontenc}
\usepackage[english]{babel}

% Page geometry. Medical journals typically want generous margins for
% reviewer comments. We use 1 inch margins which is standard.
\usepackage[margin=1in]{geometry}

% Mathematical typesetting. We need amsmath for aligned equations
% and amssymb for special mathematical symbols used in pharmacology.
\usepackage{amsmath}
\usepackage{amssymb}
\usepackage{amsfonts}

% Figure handling. The graphicx package lets us include PNG and PDF
% figures. We also load the float package for better figure placement
% control and the caption package for journal-style captions.
\usepackage{graphicx}
\usepackage{float}
\usepackage{caption}
\usepackage{subcaption}

% Table formatting. The booktabs package produces publication-quality
% tables with proper rules. We also load tabularx for tables that
% adjust their width to fit content.
\usepackage{booktabs}
\usepackage{tabularx}
\usepackage{multirow}

% Citation management. We use natbib with the apalike style which is
% common in medical journals. The actual journal will likely require
% a different style, but apalike is a reasonable default for drafting.
\usepackage{natbib}
\bibliographystyle{apalike}

% Hyperlinks within the document. The hyperref package makes the
% table of contents clickable and properly links cross-references.
\usepackage[hidelinks]{hyperref}

% Code listings for showing snippets of methodology
\usepackage{listings}
\lstset{
    basicstyle=\ttfamily\small,
    breaklines=true,
    columns=flexible
}

% Line numbering for the draft so reviewers can reference specific
% lines. Comment this out for the final submission.
\usepackage{lineno}
\linenumbers

% Color for highlighting placeholder text during drafting
\usepackage{color}
\definecolor{placeholdergray}{gray}{0.4}

% Custom command for placeholder text. Using a command lets us easily
% find and replace all placeholders later by changing this one definition.
% In the final version we will redefine this to do nothing or remove it.
\newcommand{\placeholder}[1]{\textcolor{placeholdergray}{\textit{[#1]}}}

% =====================================================================
% TITLE PAGE INFORMATION
% =====================================================================

\title{
    A Hybrid Mechanistic-Machine Learning Framework for ICU Sedation
    Prediction Through Patient-Specific Digital Twins
}

\author{
    Christopher Morris\thanks{Corresponding author. Email: \texttt{[your email here]}} \\
    Department of Computer and Information Science and Engineering \\
    University of Florida \\
    Gainesville, FL 32611, USA
}

\date{\today}

% =====================================================================
% BEGIN DOCUMENT BODY
% =====================================================================

\begin{document}

\maketitle

% =====================================================================
% ABSTRACT
% =====================================================================
%
% Medical journals typically allow 250 to 300 words for the abstract.
% Most prefer a structured abstract with explicit Background, Methods,
% Results, and Conclusions sections. We use that structure here.
%
% The abstract is the most important paragraph of your paper because
% it determines whether reviewers and readers will engage with the
% rest of the work. It should clearly state the problem, your approach,
% your key findings, and the significance of the work.
%
% Write the abstract LAST after the rest of the paper is complete.
% Use this template as a placeholder during drafting.
% =====================================================================

\begin{abstract}

\noindent\textbf{Background.}
ICU sedation management requires balancing patient comfort with safety,
but population based pharmacokinetic models fail to capture individual
variation. Pure machine learning approaches lack the interpretability
and physiological grounding required for clinical adoption.

\noindent\textbf{Methods.}
We developed a hybrid framework combining the Eleveld 2018 propofol
pharmacokinetic-pharmacodynamic model with multiple machine learning
algorithms. Patient specific effect site equilibrium concentrations
were calibrated individually using retrospective MIMIC-IV version 3.1
data from \placeholder{N patients with M observations}. Mechanistic
features including effect site concentration, plasma concentration,
and patient calibrated EC50 were augmented to clinical features and
tested across five algorithm families. We performed ablation studies
to identify which feature types contribute most to predictive
performance for each algorithm family.

\noindent\textbf{Results.}
Mechanistic feature augmentation produced consistent improvements in
mean absolute error across \placeholder{specify number} of five
tested algorithm families. The full hybrid model achieved
\placeholder{insert MAE with 95\% CI}. Ablation analysis revealed that
\placeholder{XGBoost / LSTM / both} relied primarily on
\placeholder{static / temporal / both types of} mechanistic features,
suggesting that mechanistic feature type should be matched to
algorithm architecture.

\noindent\textbf{Conclusions.}
Hybrid mechanistic-ML frameworks improve ICU sedation prediction across
multiple algorithm families. The differential pattern by which different
algorithms benefit from mechanistic features provides a roadmap for
designing physiologically grounded ML systems. Patient specific digital
twins constructed from population pharmacokinetic models offer a
clinically interpretable path to personalized sedation management.

\end{abstract}

\noindent\textbf{Keywords:} digital twin, pharmacokinetic-pharmacodynamic
modeling, propofol, ICU sedation, machine learning, hybrid modeling,
MIMIC-IV

% =====================================================================
% INTRODUCTION
% =====================================================================
%
% The introduction motivates the research question by describing what
% is known and what gap exists. It typically follows a funnel structure
% where you start broad with the clinical problem, narrow to existing
% approaches and their limitations, and end with your specific
% contribution. Aim for 800 to 1200 words.
% =====================================================================

\section{Introduction}

\subsection{Clinical Problem}

Sedation management in the intensive care unit presents a fundamental
clinical challenge that affects patient outcomes, length of stay, and
healthcare costs. Mechanically ventilated patients require pharmacological
sedation to tolerate the discomfort of intubation and to facilitate care
delivery, but the optimal sedation level depends on individual patient
factors that vary substantially across the population. Insufficient
sedation can lead to patient agitation, accidental device removal,
ventilator dysynchrony, and increased oxygen consumption. Excessive
sedation prolongs mechanical ventilation, increases the risk of delirium,
contributes to ICU acquired weakness, and extends ICU and hospital stays
\placeholder{cite recent ICU sedation outcomes literature}.

The Sedation Agitation Scale provides a standardized clinical instrument
for quantifying sedation depth on a seven point ordinal scale ranging
from unarousable to dangerously agitated. Clinicians typically aim to
maintain patients at SAS scores of three to four, which represent
appropriate sedation for mechanical ventilation while preserving the
ability to assess neurological status and engage in physical therapy.
Despite the availability of validated assessment tools, achieving and
maintaining target sedation levels remains difficult because of substantial
inter individual variation in pharmacokinetic and pharmacodynamic
responses to sedative medications \placeholder{cite Eleveld 2018 and
related personalization literature}.

\subsection{Limitations of Current Approaches}

Two complementary approaches have been applied to the prediction of
sedation responses, each with distinct strengths and limitations.
Mechanistic pharmacokinetic and pharmacodynamic models, exemplified
by the Eleveld 2018 propofol model, encode the physiological processes
governing drug distribution and effect using systems of differential
equations. These models offer interpretability and physiological
grounding, but their reliance on population averaged parameters limits
their ability to capture individual variation. Two patients receiving
identical propofol infusions will be predicted to reach similar effect
site concentrations under the population model, even though their actual
sedation responses may differ substantially due to factors not captured
in the model parameters.

Pure machine learning approaches, including gradient boosted decision
trees and recurrent neural networks, learn predictive patterns directly
from clinical data without imposing structural assumptions about the
underlying physiology. These approaches can capture individual variation
when sufficient training data is available, but they suffer from limited
interpretability and may extrapolate poorly outside the conditions
represented in their training data. Clinicians often distrust black box
predictions for high stakes decisions, and regulatory bodies have
expressed reservations about deploying purely data driven models in
clinical settings without clear physiological justification.

The limitations of each approach in isolation suggest that a hybrid
framework combining mechanistic structure with data driven learning
could outperform either component alone. This intuition has been
explored in cardiovascular digital twin research, where Corral-Acero
and colleagues described a synergy between mechanistic and statistical
models that improves both clinical interpretability and predictive
accuracy \placeholder{cite Corral-Acero 2020}. However, similar hybrid
approaches have not been systematically evaluated for ICU sedation
prediction, and the question of how different machine learning
architectures interact with mechanistic features remains open.

\subsection{Research Contributions}

This paper makes three contributions to the literature on physiologically
grounded machine learning for clinical prediction. First, we demonstrate
that augmenting machine learning models with mechanistic features
derived from the Eleveld 2018 propofol pharmacokinetic pharmacodynamic
model improves prediction of Sedation Agitation Scale scores in ICU
patients. We test this hypothesis across five algorithm families
including tree based methods, support vector machines, feedforward
neural networks, and recurrent neural networks, providing evidence
that the benefit of mechanistic augmentation is not specific to any
single algorithm family.

Second, we identify that different algorithm families benefit from
different categories of mechanistic features through systematic
ablation analysis. \placeholder{Insert specific finding once results
are finalized, e.g., recurrent neural networks benefit most from
temporal effect site concentration sequences while tree based models
benefit most from static patient calibrated parameters}. This finding
provides methodological guidance for future researchers designing
hybrid mechanistic machine learning systems by suggesting that
mechanistic feature type should be matched to algorithm architecture.

Third, we develop a patient specific digital twin approach that
calibrates the EC50 parameter of the Eleveld model individually for
each patient using their early ICU course, then uses this personalized
parameter as input to machine learning models predicting subsequent
sedation responses. This approach combines the interpretability
advantages of mechanistic modeling with the predictive flexibility
of machine learning, providing a path toward clinically useful
personalized sedation management.

% =====================================================================
% METHODS
% =====================================================================
%
% The methods section needs to be detailed enough that another researcher
% could reproduce your work. For medical journals, this typically means
% being explicit about every methodological choice. The section is
% organized in a specific order: first we describe the data, then the
% mechanistic model that processes the data, then the machine learning
% algorithms that consume the mechanistic features, then how everything
% is evaluated. This ordering follows the data flow through the analysis
% pipeline.
% =====================================================================

\section{Methods}

\subsection{Data Source and Cohort}

\subsubsection{Database Access}

This study uses the Medical Information Mart for Intensive Care version
4 release 3.1 database, a publicly available critical care database
maintained at the Massachusetts Institute of Technology Laboratory for
Computational Physiology \placeholder{cite Johnson MIMIC-IV 2023}. The
database contains de-identified electronic health records from
approximately 300,000 patients admitted to Beth Israel Deaconess Medical
Center between 2008 and 2019. Access was obtained through credentialed
PhysioNet membership including completion of required ethics training
and signed data use agreement. Data extraction was performed via Google
BigQuery against the physionet-data project.

\subsubsection{Cohort Selection}

The patient cohort for this analysis consisted of adult ICU patients
who received propofol infusions during their stay and had sufficient
Sedation Agitation Scale assessments to support our temporal modeling
approach. Specific inclusion criteria were patient age of at least
18 years at admission, receipt of at least one propofol infusion event
during the ICU stay, and at least 11 SAS observations during the stay
to enable formation of temporal sequences with held out prediction
targets. \placeholder{Final cohort size: N patients meeting these
criteria}.

\subsubsection{Variables Extracted}

For each patient meeting inclusion criteria, we extracted demographic
variables including age, gender, weight, and height. We extracted all
propofol infusion events with start times, end times, infusion rates,
and units of measurement. We extracted all SAS observations recorded
during the ICU stay with associated timestamps. The complete SQL queries
used for data extraction are provided in the supplementary materials
and in our public code repository.

\subsection{Mechanistic Pharmacokinetic-Pharmacodynamic Model}

\subsubsection{Three-Compartment Pharmacokinetic Structure}

We use the propofol pharmacokinetic pharmacodynamic model published by
Eleveld and colleagues in 2018 \placeholder{cite Eleveld 2018}. The
pharmacokinetic component describes drug distribution through three
compartments representing the central blood compartment, a rapidly
equilibrating peripheral compartment, and a slowly equilibrating
peripheral compartment. The dynamics are governed by the system of
ordinary differential equations

\begin{align}
\frac{dA_1}{dt} &= R(t) - \frac{CL}{V_1}A_1 - \frac{Q_2}{V_1}A_1
                + \frac{Q_2}{V_2}A_2 - \frac{Q_3}{V_1}A_1 + \frac{Q_3}{V_3}A_3 \\
\frac{dA_2}{dt} &= \frac{Q_2}{V_1}A_1 - \frac{Q_2}{V_2}A_2 \\
\frac{dA_3}{dt} &= \frac{Q_3}{V_1}A_1 - \frac{Q_3}{V_3}A_3
\end{align}

where $A_1$, $A_2$, and $A_3$ represent drug amounts in the central,
rapid peripheral, and slow peripheral compartments respectively, $V_1$,
$V_2$, and $V_3$ are the corresponding compartment volumes, $CL$ is
metabolic clearance, $Q_2$ and $Q_3$ are inter-compartmental clearance
rates, and $R(t)$ is the time varying drug infusion rate.

The effect site concentration $C_e$ tracks the concentration affecting
the brain through the equation

\begin{equation}
\frac{dC_e}{dt} = k_{e0}\left(\frac{A_1}{V_1} - C_e\right)
\end{equation}

where $k_{e0}$ is the effect site equilibration rate constant.

\subsubsection{Sigmoid Emax Pharmacodynamic Response}

The pharmacodynamic component links effect site concentration to
predicted sedation level through a sigmoid Emax relationship

\begin{equation}
\widehat{SAS}(C_e) = E_0 - E_{max} \cdot \frac{C_e^\gamma}{EC_{50}^\gamma + C_e^\gamma}
\end{equation}

where $E_0$ is the baseline SAS in the absence of drug, $E_{max}$ is
the maximum possible drug effect on the SAS scale, $\gamma$ is the
Hill coefficient controlling the steepness of the dose response curve,
and $EC_{50}$ is the effect site concentration producing half maximum
effect. The $EC_{50}$ parameter represents individual drug sensitivity
and is the primary target of patient specific calibration in our
approach.

\subsubsection{Population Parameter Values}

We use the population parameter values reported by Eleveld and colleagues
which represent typical values for the average adult patient. These
values are summarized in Table~\ref{tab:eleveld_params}.

\begin{table}[H]
\centering
\caption{Eleveld 2018 population pharmacokinetic-pharmacodynamic parameter values used in our analysis.}
\label{tab:eleveld_params}
\begin{tabular}{lcl}
\toprule
Parameter & Value & Description \\
\midrule
$V_1$ & 6.28 L & Central compartment volume \\
$V_2$ & 25.5 L & Rapid peripheral compartment volume \\
$V_3$ & 273 L & Slow peripheral compartment volume \\
$CL$ & 1.79 L/min & Metabolic clearance \\
$Q_2$ & 1.75 L/min & Rapid intercompartmental clearance \\
$Q_3$ & 1.11 L/min & Slow intercompartmental clearance \\
$k_{e0}$ & 0.146 /min & Effect site equilibration rate \\
$EC_{50}^{pop}$ & 3.08 \textmu g/mL & Population EC50 \\
$\gamma$ & 1.47 & Hill coefficient \\
$E_0$ & 7 & Baseline SAS without drug \\
$E_{max}$ & 6 & Maximum drug effect on SAS \\
\bottomrule
\end{tabular}
\end{table}

\subsubsection{Patient-Specific EC50 Calibration}

The core innovation of our digital twin approach is the calibration
of $EC_{50}$ to each individual patient. For each patient we use the
first half of their SAS observations as a calibration set, finding the
$EC_{50}$ value that minimizes the mean squared error between
mechanistic predictions and observed SAS scores at the calibration
time points. The calibration is performed via bounded scalar
optimization over the range from 0.5 to 10.0 micrograms per milliliter
using Brent's method.

The calibrated $EC_{50}$ is used to compute mechanistic predictions
for the held out second half of the patient's observations, which
serve as the prediction targets for our machine learning models. This
calibration approach mimics how the system would function clinically,
where early ICU observations would be used to personalize the model
before generating predictions for subsequent care decisions.

\subsection{Feature Engineering}

We construct feature vectors for each prediction instance representing
one moment where we want to predict the next SAS score. The feature
vector includes ten clinical features derived from patient demographics
and recent SAS history, plus five static mechanistic features representing
the pharmacokinetic state at prediction time and including the patient
calibrated EC50. For temporal models, we additionally provide a sequence
of effect site concentrations at the past ten observation time points.

The complete list of features and their derivation is provided in
Table~\ref{tab:features}.

\begin{table}[H]
\centering
\caption{Features used in machine learning models, organized by category.}
\label{tab:features}
\begin{tabularx}{\textwidth}{lXl}
\toprule
Feature & Description & Category \\
\midrule
age & Patient age in years at ICU admission & Clinical \\
gender & Binary encoded (1=Male, 0=Female) & Clinical \\
weight & Body weight in kilograms & Clinical \\
height & Height in centimeters & Clinical \\
sas\_last & Most recent SAS observation before prediction time & Clinical \\
sas\_mean & Mean of last 10 SAS observations & Clinical \\
sas\_std & Standard deviation of last 10 SAS observations & Clinical \\
sas\_trend & Change between earliest and latest in 10 observation window & Clinical \\
propofol\_rate & Current infusion rate at prediction time & Clinical \\
time\_hours & Time since admission in hours & Clinical \\
\midrule
$C_e$ & Effect site concentration at prediction time & Static mechanistic \\
$C_p$ & Plasma concentration at prediction time & Static mechanistic \\
$EC_{50}^{cal}$ & Patient calibrated EC50 value & Static mechanistic \\
mech\_pred\_pop & Mechanistic SAS prediction using population EC50 & Static mechanistic \\
mech\_pred\_calib & Mechanistic SAS prediction using calibrated EC50 & Static mechanistic \\
\midrule
$C_e$ sequence & Last 10 effect site concentrations as a sequence & Temporal mechanistic \\
SAS sequence & Last 10 SAS observations as a sequence & Temporal mechanistic \\
\bottomrule
\end{tabularx}
\end{table}

\subsection{Machine Learning Models}

\subsubsection{Algorithm Selection}

We evaluated five machine learning algorithms representing distinct
algorithm families. XGBoost \placeholder{cite Chen Guestrin 2016}
represents gradient boosted decision trees. Random Forest
\placeholder{cite Breiman 2001} represents bagged tree ensembles.
Support Vector Regression with radial basis function kernel
\placeholder{cite Smola Schölkopf 2004} represents kernel based
methods. Multi Layer Perceptron represents feedforward neural networks.
Long Short Term Memory \placeholder{cite Hochreiter Schmidhuber 1997}
represents recurrent neural networks designed for sequential data.

These five algorithms span the major categories of machine learning
approaches used in medical informatics, providing breadth in our
evaluation of how mechanistic feature augmentation interacts with
different algorithm families.

\subsubsection{Hyperparameter Specifications}

\placeholder{Insert detailed hyperparameter specifications for each
algorithm here, drawing from the transparency checklist. Include
justifications for each choice.}

\subsection{Ablation Study Design}

To identify which feature categories contribute most to predictive
performance for each algorithm family, we conducted a systematic
ablation study testing five feature configurations. The Full Hybrid
configuration uses all features and serves as the reference point.
The No Temporal configuration removes the temporal mechanistic features
while retaining static features. The No Static Mechanistic configuration
removes the static mechanistic features while retaining temporal features.
The Pure ML configuration removes all mechanistic features, leaving
only clinical features. The Only Mechanistic configuration removes
clinical features, retaining only mechanistic features.

For each configuration, we trained XGBoost once due to its deterministic
behavior given a fixed random seed. We trained LSTM three times with
different random seeds to assess result stability across initializations
because neural network training is stochastic. We report mean and
standard deviation of test set performance across the three LSTM seeds
for each configuration.

\subsection{Evaluation Methodology}

\subsubsection{Train Test Split}

We use an 80 to 20 patient level train test split. The split is patient
level rather than observation level to prevent data leakage, ensuring
that the same patient never appears in both training and test sets.
This patient level splitting provides a more honest estimate of how
the model would generalize to new patients in clinical deployment.

\subsubsection{Performance Metrics}

The primary performance metric is mean absolute error in SAS units,
which directly quantifies the typical magnitude of prediction errors
in clinically interpretable units. We additionally report root mean
squared error which penalizes larger errors more heavily, and
coefficient of determination R squared which represents the fraction
of variance explained by the model.

% =====================================================================
% RESULTS
% =====================================================================
%
% The results section presents what was found, organized to support
% the claims made in the introduction. We follow the order of the
% hypotheses we are testing, with the primary result first.
%
% Critical: do not interpret results in this section. Save interpretation
% for the discussion. The results section should be a factual report
% of what the data show.
% =====================================================================

\section{Results}

\subsection{Cohort Characteristics}

\placeholder{Insert Table showing patient demographics including total
N, gender breakdown, age distribution, BMI distribution, etc. Use the
standard medical journal Table 1 format.}

\subsection{Pharmacokinetic Model Performance}

\placeholder{Insert results showing the performance of the population
Eleveld model alone, then the patient calibrated version, demonstrating
that calibration improves predictions.}

\subsection{Multi-Algorithm Hybrid Comparison}

We compared each of the five algorithm families in pure machine learning
configuration using only clinical features against the same algorithms
augmented with mechanistic features in the hybrid configuration. The
results are summarized in Table~\ref{tab:multialg}.

\begin{table}[H]
\centering
\caption{Performance comparison between pure machine learning and hybrid
configurations across five algorithm families. Results show mean absolute
error on the test set with 95\% confidence intervals.}
\label{tab:multialg}
\begin{tabular}{lccc}
\toprule
Algorithm & Pure ML MAE & Hybrid MAE & Improvement \\
\midrule
XGBoost & \placeholder{XX} & \placeholder{XX} & \placeholder{X.X\%} \\
Random Forest & \placeholder{XX} & \placeholder{XX} & \placeholder{X.X\%} \\
SVM (RBF) & \placeholder{XX} & \placeholder{XX} & \placeholder{X.X\%} \\
MLP & \placeholder{XX} & \placeholder{XX} & \placeholder{X.X\%} \\
LSTM & \placeholder{XX} & \placeholder{XX} & \placeholder{X.X\%} \\
\bottomrule
\end{tabular}
\end{table}

\placeholder{Insert Figure 1 here showing the bar chart visualization
of these results.}

\subsection{Ablation Study Results}

The ablation study results are summarized in Table~\ref{tab:ablation}
and visualized in Figure~\placeholder{X}.

\begin{table}[H]
\centering
\caption{Ablation study results showing performance degradation from
the Full Hybrid baseline for each feature configuration. LSTM results
report mean and standard deviation across three random seeds.}
\label{tab:ablation}
\begin{tabular}{lcccc}
\toprule
& \multicolumn{2}{c}{XGBoost} & \multicolumn{2}{c}{LSTM} \\
\cmidrule(lr){2-3}\cmidrule(lr){4-5}
Configuration & MAE & Degradation & MAE & Degradation \\
\midrule
Full Hybrid & \placeholder{XX} & --- & \placeholder{XX$\pm$X} & --- \\
No Temporal & \placeholder{XX} & \placeholder{X.X\%} & \placeholder{XX$\pm$X} & \placeholder{X.X$\pm$X.X\%} \\
No Static Mechanistic & \placeholder{XX} & \placeholder{X.X\%} & \placeholder{XX$\pm$X} & \placeholder{X.X$\pm$X.X\%} \\
Pure ML & \placeholder{XX} & \placeholder{X.X\%} & \placeholder{XX$\pm$X} & \placeholder{X.X$\pm$X.X\%} \\
Only Mechanistic & \placeholder{XX} & \placeholder{X.X\%} & \placeholder{XX$\pm$X} & \placeholder{X.X$\pm$X.X\%} \\
\bottomrule
\end{tabular}
\end{table}

\placeholder{Describe the key finding of the ablation study in narrative
form, focusing on how XGBoost and LSTM differ in their response to
removing different feature types. This is the core empirical finding
of the paper.}

% =====================================================================
% DISCUSSION
% =====================================================================
%
% The discussion interprets the results in context. It typically has
% several subsections covering principal findings, comparison with
% existing literature, clinical implications, limitations, and future
% directions. Aim for 1500 to 2000 words.
% =====================================================================

\section{Discussion}

\subsection{Principal Findings}

\placeholder{Summarize the key findings in narrative form. State what
the experiments showed, focusing on the headline results. Reserve
about three paragraphs for this subsection.}

\subsection{Relationship to Existing Literature}

The hybrid mechanistic machine learning approach we develop builds
on the model synergy concept articulated by Corral-Acero and colleagues
in cardiovascular digital twin research \placeholder{cite Corral-Acero
2020}. They proposed that mechanistic and statistical models can be
combined to produce predictions that are both interpretable and
accurate, with mechanistic models providing physiological grounding
and statistical models providing flexibility. Our work applies this
general principle to a specific clinical problem in ICU sedation
management and provides empirical evidence that the synergy operates
across multiple machine learning algorithm families.

\placeholder{Add comparisons with other relevant work on hybrid
approaches and on machine learning for sedation prediction.}

\subsection{Methodological Implications}

\placeholder{Discuss what the ablation study results imply for future
work designing hybrid mechanistic machine learning systems. The key
insight is that algorithm architecture should be matched to mechanistic
feature type. Discuss this in detail.}

\subsection{Clinical Implications}

\placeholder{Discuss what the findings mean for clinical practice.
Patient specific dosing, reduced sedation errors, potential for
personalized critical care. Be careful not to overclaim, since this
is retrospective analysis on a single dataset.}

\subsection{Limitations}

\placeholder{Honest discussion of limitations including retrospective
single center data, the SAS to BIS validation gap in the Eleveld model,
remaining LSTM training instability, and limited cohort diversity.}

\subsection{Future Directions}

\placeholder{Discuss the natural next steps including prospective
validation, multi center external validation, real time deployment,
integration with electronic health records, and the planned reinforcement
learning extension for personalized sedation management.}

% =====================================================================
% CONCLUSION
% =====================================================================
%
% A brief paragraph summarizing the key contributions and significance.
% Should not introduce new ideas. About 200 words.
% =====================================================================

\section{Conclusion}

\placeholder{Write the conclusion last, after all other sections are
complete. Should briefly restate the contribution, summarize key
findings, and indicate the broader significance.}

% =====================================================================
% ACKNOWLEDGMENTS
% =====================================================================

\section*{Acknowledgments}

\placeholder{Include acknowledgments to advisor, funding sources,
and any other contributors. Standard text typically thanks PhysioNet
for providing access to MIMIC-IV.}

% =====================================================================
% DATA AND CODE AVAILABILITY
% =====================================================================
%
% Most medical journals now require this section. Be specific about
% where code can be found and how data access works.
% =====================================================================

\section*{Data and Code Availability}

The MIMIC-IV database used in this study is publicly available to
credentialed researchers through PhysioNet at
\url{https://physionet.org/content/mimiciv/3.1/}. Access requires
completion of required ethics training and signed data use agreement.
The patient cohort used in this analysis can be reproduced by running
the SQL queries provided in our public code repository against
MIMIC-IV version 3.1.

All analysis code is available at \placeholder{insert repository URL}
under the \placeholder{insert license name} license. The repository
includes the data extraction queries, feature processing code, model
training scripts, statistical analysis code, and documentation
necessary to reproduce all results reported in this paper.

% =====================================================================
% AUTHOR CONTRIBUTIONS
% =====================================================================
%
% Some journals require this. Standard format below.
% =====================================================================

\section*{Author Contributions}

\placeholder{Standard author contributions statement. C.M. designed
the research, conducted the analyses, and drafted the manuscript.
[Advisor name] supervised the research and contributed to the manuscript.}

% =====================================================================
% CONFLICTS OF INTEREST
% =====================================================================

\section*{Conflicts of Interest}

The authors declare no conflicts of interest.

% =====================================================================
% REFERENCES
% =====================================================================
%
% References will be managed through a .bib file. The bibliography
% style is set in the preamble. For now we use placeholder citations
% in the text that we will replace with proper references.
% =====================================================================

\bibliography{references}

\end{document}
